% !TeX program = xelatex
\documentclass[11pt,a4paper]{article}
\usepackage[margin=23mm]{geometry}
\usepackage{fontspec}
\setmainfont{Latin Modern Roman}
\usepackage{amsmath,amssymb,mathtools}
\usepackage{xcolor,enumitem,booktabs,longtable,array,xurl,hyperref}
\definecolor{muted}{HTML}{53636C}
\hypersetup{colorlinks=true,urlcolor=blue,linkcolor=blue}
\setlength{\parindent}{0pt}
\setlength{\parskip}{5pt}
\newcommand{\R}{\mathbb R}
\newcommand{\E}{\mathbb E}
\newcommand{\N}{\mathbb N}
\newcommand{\ind}{\mathbf 1}
\newcommand{\Var}{\operatorname{Var}}
\newcommand{\Cov}{\operatorname{Cov}}
\newcommand{\supp}{\operatorname{supp}}
\DeclareMathOperator*{\argmin}{arg\,min}
\DeclareMathOperator*{\argmax}{arg\,max}
\newcommand{\entrymeta}[3]{{\sffamily\small\color{muted}\textbf{\texttt{#1}}\quad|\quad #2\par #3\par}}
\newcommand{\Problemref}[1]{\texttt{#1}}
\newcommand{\JELref}[2]{\texttt{#2} (JEL \texttt{#1})}
\begin{document}
\section*{Home-school investment complementarities}
\textbf{Problem ID:} \texttt{I21-615}\quad \textbf{JEL:} \texttt{I21}\par
% BEGIN ECONOMICS STATEMENT
\label{problem:OP-0451}
{\sffamily\small\color{muted}Original subject group: Education, families and demography\par}
\label{definitions:problem:30007733}
\textbf{Audit classification:} Published research agenda. Combining utility and monetary program costs requires a declared normalization; the welfare specification below leaves that choice open.
\subsection*{Definitions and precise research target}
\textit{Editorial formalization.} Consider four-year-old children in disadvantaged households in one urban school district. Let \(p\in\{0,1\}\) be a randomized offer of a specified preschool curriculum and \(h\in\{0,1\}\) a randomized offer of a home-visiting parenting program. Define \(S_i(p,h)\) as independently assessed school readiness one year later, and \(V_i(p,h)\) as discounted family welfare over five years net of participation time and program resource costs. The interaction targets are
\[\gamma_S=E[S_i(1,1)-S_i(1,0)-S_i(0,1)+S_i(0,0)],\]
and \(\gamma_V\), defined by replacing readiness with welfare. Identify both quantities using factorial offers, with consistency and either negligible household spillovers or a stated neighborhood exposure design. Positive \(\gamma_S\) establishes complementarity of these offers on this outcome scale, not necessarily technological complementarity of realized educational inputs. Collect attendance, parental time, private expenditure, and curriculum implementation to evaluate that distinction. Report interaction uncertainty and baseline heterogeneity instead of inferring complementarity from separate positive average effects. The review explicitly calls for combined preschool and parenting evidence and studying parental choices. Its general agenda is adapted here to defined programs and welfare accounting; no sign restriction on either interaction is assumed.

Let \(p,h\) denote offered programs, not realized attendance/time inputs. Readiness \(S\) is an independently scaled bounded score. Define welfare \(V_i=E\sum_{t=0}^{4}\beta^tu_i(c_{it},\text{childhealth}_{it},\text{leisure}_{it})\), with program costs charged through a declared budget; subtracting money from utility requires a specified quasilinear normalization. Alternatively use a baseline-money equivalent of this expected utility and subtract resource costs only once. All four offer combinations have positive randomization probability. Positive \(\gamma_S\) means strict supermodularity of the four mean offer outcomes; weak supermodularity also permits \(\gamma_S=0\). This property is scale-specific and need not imply \(\partial^2F/\partial p\partial h>0\) for any technology using actual educational inputs. \(\gamma_V\) additionally depends on preferences, participation costs, and financing. All expectations use a stated baseline population and probability law. Identification means every admissible data-generating model with the same observed-data law yields the same displayed target; otherwise report the identified set. Exact equations, horizons and assumptions are editorial, rather than source-given canonical conjectures.

The welfare expression is a specification template: choosing utility accounting or a baseline-money equivalent, family aggregation and population weights, cardinal normalization, discounting, and the information and probability law for its inner expectation remains necessary to define a particular \(\gamma_V\). Any monetary equivalent must specify the transfer and price convention and the domain on which its defining equation has a solution. The accounting perspective, financing incidence and allocation of resource costs must be declared, identifying where participation time and program costs enter so that each is counted once. No particular selection among these choices is imposed here.

The identification task also requires a declared observed-data law: specify the one-year readiness and five-year welfare observations, which welfare primitives are known or separately identified, and the treatment of missing outcomes. Randomized offers alone do not identify unknown welfare preferences or unobserved welfare components. State assignment exchangeability and support relative to any baseline randomization strata and the target population; identification within randomized participants does not by itself establish generalization to all eligible district children. Retain the identified-set criterion above when point identification is unavailable. The district, baseline eligibility and sampling frame, program versions and durations, no-offer services, and common readiness scoring rule are study primitives to be declared before empirical application.

The interference alternative remains open. For exact identification, negligible spillovers cannot substitute for a defined no-interference assumption without an approximation criterion. A neighborhood or classroom exposure formulation must specify exposure-indexed potential outcomes, the exposure held fixed or averaged over in each contrast, and support for the required joint own-offer and exposure regimes. Own-offer arm positivity alone does not provide that support. Neither an exposure regime nor a tolerance for approximate interference is selected here.
\subsection*{Variants and assumption changes}
\begin{enumerate}
\item \textbf{Editorial actual-input variant.} Replace offers by instructional hours and parental time in a declared production function. Input cross-effects require valid excluded variation and support.
\item \textbf{Editorial caregiver-stress variant.} Add randomized caregiver mental-health support, creating eight arms. Define third-order interactions and account for parent/child welfare costs once.
\end{enumerate}
\subsection*{References and source audit}
\textbf{1.} Explicit agenda on long-run curricula, follow-on investments, and combined preschool/parenting programs. Locator: BFI Working Paper 2022-58; Section 7, printed pp.91–93, particularly p.93. Full text or primary publication page inspected. URL: \url{https://bfi.uchicago.edu/wp-content/uploads/2022/05/BFI_WP_2022-58.pdf}.
\begin{enumerate}\raggedright
\item\hypertarget{definitions:source:30007733:1}{} Greg J. Duncan; Ariel Kalil; Magne Mogstad; Mari Rege. 2022. \emph{Investing in Early Childhood Development in Preschool and at Home}. BFI Working Paper 2022-58; Section 7, printed pp.91–93, particularly p.93.
\url{https://bfi.uchicago.edu/wp-content/uploads/2022/05/BFI_WP_2022-58.pdf}.
DOI: \url{https://doi.org/10.3386/w29985}.
\end{enumerate}

\subsection*{Common conventions: Source mathematical conventions}

These conventions apply to each entry unless explicitly replaced.

\textbf{Models and identification.} A primitive model \(m\) specifies all probability laws, feasible choices, information, equilibrium and measurement rules used by its target. The admissible class \(\mathcal M\) is fixed independently of outcomes. If \(P_O(m)\) is its observable probability law and \(\tau(m)\) the target, its identified set at observable law \(P\) is
\[
 \mathcal I_\tau(P)=\{\tau(m):m\in\mathcal M,\ P_O(m)=P\}.
\]
Point identification means that this set is a singleton. An empty set signals incompatibility of data and assumptions. Coordinate bounds need not describe the joint set. Bounds are sharp only if every retained target is attainable or arbitrarily approachable by compatible models. Finite bounds require explicit restrictions on unobserved quantities; unrestricted confounding, hidden wealth, extrapolation or arbitrary equilibrium selection can yield uninformative sets.

\textbf{Causal interventions.} A potential outcome \(Y_i(p)\) is the outcome under a complete policy regime \(p\), including timing, financing, information and market spillovers. The target population and units are fixed before treatment. With interference, use the complete assignment vector or a justified exposure mapping; individual no-interference notation alone cannot identify an economy-wide intervention. A difference of mean potential outcomes does not require the cross-world joint distribution. Conditioning on post-treatment migration, survival, enrollment or employment changes the estimand and needs separate justification.

\textbf{Statistical statements.} An estimator, confidence set, forecast or test specifies a sampling law, model class and loss/coverage criterion. Uniform \((1-\alpha)\)-coverage means \(\inf_{m\in\mathcal M}\Pr_m\{\tau(m)\in C_n\}\geq1-\alpha\), or an explicitly asymptotic version. The whole parameter space trivially has coverage, so informativeness requires a stated width, risk or power objective. A forecasting improvement is relative to a named benchmark, information set and evaluation population.

\textbf{Equilibrium and optimization.} An equilibrium comprises feasible plans, optimality, beliefs where needed, budget constraints and all market/resource clearing. The primitives or a stated benchmark define each correspondence \(\mathcal E(p;m)\). Nonemptiness is assumed or proved, never inferred just from compact policy sets. A selected-equilibrium target uses a declared selection rule; a robust target quantifies over all permitted equilibria. Use suprema or infima unless attainment is established. An \(\varepsilon\)-optimal policy is feasible and within \(\varepsilon\) of the finite optimum value. Report an empty feasible set separately.

\textbf{Welfare and accounting.} Normative utility, interpersonal normalization, social weights, numeraire, discounting and the use of public receipts are primitives. Behavioral utility need not coincide with the welfare criterion. Household welfare includes ownership income and rents once; adding profits again to compensating variation generally double-counts them. A monetary equivalent is defined by an explicit transfer and price convention, with monotonicity and range conditions. Average effects, mechanism contributions, transfers and welfare losses are different targets. Infinite sums require convergence and dynamic feasible plans require terminal or no-Ponzi conditions.

\textbf{Source versions.} A working-paper date, revision date and journal publication year are distinguished. A DOI for a journal version is not automatically the DOI of an earlier manuscript. Locators refer to the stated version and printed pages where specified. A metadata-only check does not verify a claim about a full-text conclusion. Every entry's source subsection narrows the provenance claim accordingly.
% END ECONOMICS STATEMENT
\end{document}
